\section{Introduction}\label{sec:introduction}
The Singer conjecture predicts that the $L^2$-homology of the universal
cover of a closed aspherical manifold is concentrated in its middle
dimension. In dimension four, Poincar\'e duality reduces the problem to
vanishing in degree one. The conclusion forces the Euler characteristic
to be nonnegative. We prove this vanishing for closed aspherical
topological four-manifolds
and, more generally, for finite aspherical integral Poincar\'e complexes
of formal dimension four.

A connected CW complex or manifold is \emph{aspherical} if its universal
cover is contractible. Throughout, a closed manifold is compact and has
no boundary. A \emph{finite integral Poincar\'e complex of formal dimension
four} means a finite connected CW complex $Q$ equipped with an orientation
character $w:\pi_1(Q)\to\{1,-1\}$ and a fundamental class
$[Q]\in H_4(Q;\mathbb Z^w)$ for which cap product gives duality in every
degree with every local coefficient system, with the orientation twist
on homology. This is absolute duality for $Q$, not duality for a
Poincar\'e pair. Equivalently, the duality condition can be tested on the
regular group-ring local system \cite[Lemma~1.1]{Wall1967}.

For a discrete group $\Gamma$, write $\mathcal N\Gamma$ for its group
von Neumann algebra and $\dim_{\mathcal N\Gamma}$ for the extended
Murray--von Neumann dimension. If $\widetilde Q\to Q$ is the universal
cover, with deck group $\Gamma=\pi_1(Q)$, set
\[
 b_p^{(2)}(\widetilde Q)
 =\dim_{\mathcal N\Gamma}H_p\bigl(
 \mathcal N\Gamma\otimes_{\mathbb Z\Gamma}C_*^s(\widetilde Q)\bigr),
\]
where $C_*^s$ denotes singular chains. This definition applies to
CW complexes and topological manifolds and agrees with the usual
cellular definition
\cite[Definition~6.50 and Lemmas~6.51--6.53]{Luck2002}.

\begin{theorem}\label{thm:main}
Let $Q$ be a finite connected aspherical integral Poincar\'e complex of
formal dimension four, with arbitrary orientation character. Then
\[
 b_p^{(2)}(\widetilde Q)=0\qquad(p\ge0,\ p\ne2).
\]
\end{theorem}

\begin{corollary}\label{cor:singer-manifolds}
Let $M$ be a closed connected aspherical topological four-manifold. Then
\[
 b_p^{(2)}(\widetilde M)=0\qquad(p\ge0,\ p\ne2).
\]
\end{corollary}

\begin{proof}
Lemma~\ref{found:model} supplies a finite aspherical integral
Poincar\'e-complex homotopy model of $M$. A homotopy equivalence lifts to
an equivariant homotopy equivalence of universal covers and preserves
the singular-chain invariants above. Theorem~\ref{thm:main} applies to
the model.
\end{proof}

Corollary~\ref{cor:singer-manifolds} establishes the Singer conjecture
in dimension four. Neither orientability nor a geometric structure on
$M$ is required. The substantive assertion is the vanishing in degree
one; Poincar\'e duality supplies the other nonmiddle degrees. In
particular, the corollary includes manifolds that admit no smooth
structure or triangulation. The finite complex used for such a manifold
is a homotopy model.

For every $Q$ in Theorem~\ref{thm:main}, the remaining invariant is
\[
 \chi(Q)=b_2^{(2)}(\widetilde Q)\ge0.
\]
Thus $\chi(Q)=0$ exactly when all its $L^2$-Betti numbers vanish.
For an oriented manifold $M$ in Corollary~\ref{cor:singer-manifolds},
the $L^2$-signature theorem also gives $\chi(M)\ge|\sigma(M)|$.
Along a nested normal tower of finite
covers with trivial subgroup intersection, $L^2$ approximation determines
the limits of normalized rational Betti numbers. Section~\ref{sec:consequences}
gives the precise statements and proofs of these consequences.

The finite-Poincar\'e-complex extension was formulated in the earlier
literature \cite[Remark~11.3]{Luck2016Approximation}.
It applies beyond manifold models: the companion
\emph{A PD4 group without an aspherical manifold model} constructs a
finite oriented aspherical integral Poincar\'e complex of formal dimension
four whose fundamental group is not that of any closed aspherical
topological four-manifold
\cite[Theorem~1.1]{OpenAIPD4Nonmanifold2026}.
Theorem~\ref{thm:main} applies to that example directly.

\subsection{History and significance}
Atiyah's $L^2$ index theorem explains the connection with Euler
characteristic.
For an elliptic operator on a closed smooth manifold, its index agrees
with the von Neumann index of the lifted operator on a normal cover
\cite[Theorem~3.8 and Section~6.4]{Atiyah1976}.
Applied to differential forms, this identifies the Euler characteristic
with the alternating sum of the $L^2$-Betti numbers. Dodziuk's
analytic--simplicial comparison and homotopy-invariance theorem
\cite[Theorems~1--2]{Dodziuk1977} connect this analytic description to
topology. The algebraic definition used here also applies to finite
complexes and topological manifolds without a smooth structure.

Dodziuk records Singer's proposal to obtain the Euler-characteristic
sign from vanishing of $L^2$-harmonic forms on universal covers
\cite[p.~396]{Dodziuk1979}. His discussion concerns nonpositive
curvature. The later aspherical-manifold formulation asks that, for a
closed aspherical $n$-manifold, the $L^2$-Betti numbers vanish whenever
$p\ne n/2$; see \cite[Conjecture~0.3]{DavisOkun2001} and
\cite[Conjecture~4.1]{KirsteinKremerLuck2025}.
For even $n=2m$, concentration would give
$(-1)^m\chi(M)=b_m^{(2)}(\widetilde M)\ge0$.
In dimension four the remaining vanishing question is $b_1^{(2)}=0$:
the degree-zero vanishing and Poincar\'e duality determine the other
nonmiddle degrees. Lott and L\"uck's work on three-manifolds
\cite[Theorem~0.1]{LottLuck1995} illustrates an earlier route through
geometric decomposition. In the closed case, their computation assumed that each prime
factor had a finite cover homotopy equivalent to a Haken, Seifert, or
hyperbolic manifold; geometrization later removed this restriction for
closed aspherical three-manifolds.

Several four-dimensional cases precede the general result.
Davis and Okun proved the vanishing for right-angled Coxeter groups
whose nerves are flag triangulations of the three-sphere
\cite[Theorem~11.1.1]{DavisOkun2001}.
Okun and Schreve established the Singer conjecture in dimensions at
most four for groups of type~$\mathrm{VF}$ admitting a hierarchy,
and for Coxeter groups whose nerves are arbitrary triangulations of
the three-sphere \cite[Theorems~4.16--4.17]{OkunSchreve2016}.
Here type~$\mathrm{VF}$ means that a finite-index subgroup has a finite
classifying space. A hierarchy, in their sense, repeatedly cuts a
contractible manifold carrying a proper cocompact group action along
invariant locally flat hypersurfaces with contractible components,
ending in compact contractible pieces. The Singer conclusion concerns
an initial manifold without boundary.
These results permit topological manifolds. Their extra input is the
specified decomposition, which is not supplied by Poincar\'e duality
alone.

Complex geometry provides another route. The closed aspherical K\"ahler
surface case follows from Gromov's restriction on K\"ahler groups with
positive first $L^2$-Betti number, as explained in
\cite[Theorem~4.1]{AlbaneseDiCerboLombardi2025}.
Albanese, Di Cerbo and Lombardi also prove Singer vanishing for closed
aspherical complex surfaces with residually finite fundamental group,
and for those outside class~$\mathrm{VII}_0^+$ without that hypothesis
\cite[Theorems~1.5 and~4.2]{AlbaneseDiCerboLombardi2025}.
Their residual-finite argument uses the Albanese map, classification
of complex surfaces and $L^2$ approximation.
Class~$\mathrm{VII}_0^+$ consists of minimal non-K\"ahler surfaces with
Kodaira dimension $-\infty$ and positive second Betti number.
The separate companion global spherical shell theorem
\cite[Theorem~1.1 and Corollary~1.2]{OpenAIGlobalSphericalShells2026}
gives these surfaces fundamental group $\mathbb Z$, excluding
asphericity because an aspherical complex surface with that group would be
homotopy equivalent to $S^1$ and have zero second Betti number.
Combined with the cited Theorem~4.2, this gives an
independent argument for the complex-surface case. It is not an input
to the proof below.

Avramidi, Okun and Schreve prove the conjecture for the
even-dimensional Gromov--Thurston cyclic branched covers arising from
a closed oriented hyperbolic manifold and two separating transverse
totally geodesic hypersurfaces \cite[Section~3, p.~2628]{AvramidiOkunSchreve2025}.
The branch locus is their intersection; the cyclic covering is
specified by algebraic intersection with a chosen half of one
hypersurface. Their cutting argument covers every positive branching degree
in this even-dimensional setting. The proof here uses the algebraic
consequences of integral Poincar\'e duality to construct a boundary and
control its topology.

Finite-energy functions have a classical role in constructing boundaries.
For graphs, the Royden compactification uses all bounded functions of
finite Dirichlet energy \cite[Section~4]{KellerLenzSchmidtWojciechowski2017}.
We select a countable invariant family and add coordinates recording
components of coordinate neighborhoods. This produces a metrizable
boundary; its convergence dynamics and graph-path properties are
proved directly. Bowditch's complete median pretree and dendritic
quotient organize the cut points
\cite{Bowditch1999Treelike,Bowditch1999Connected}.
The invariant finite-interval reduction in the all-parabolic case and
the continuous retraction onto a block require the additional arguments
given below. Our separation counts adapt Bowditch's annulus-system
method \cite[Sections~6--7]{Bowditch1998Topological}; the component
labels and the threshold uniform over all sufficiently deep paths are
proved in the present setting.
The path-pushing strategy follows the Bestvina--Mess viewpoint as
expounded by Hruska and Ruane~\cite[Sections~1 and~3]{HruskaRuane2024}.
Here annular depth, component labels, and three anchors adapt that
strategy to graphs that may have infinite valence and infinite
stabilizers; the word-hyperbolic local-connectivity theorem is not
being applied.

\subsection{The proof in outline}
Suppose the first $L^2$-Betti number is positive. A finite presentation
of $\Gamma$ gives a locally finite Cayley graph, and the corresponding
reduced $L^2$-cohomology supplies a bounded function of finite Dirichlet
energy whose behavior at infinity is nonconstant. Here finite energy
means that the sum of the squared differences across all graph edges
is finite.

We use translates of this function to build a compactification
$X=\Gamma\sqcup B$. Its boundary $B$ is a nondegenerate continuum,
meaning a compact connected metric space with more than one point.
Additional coordinates distinguish the graph
components of suitable coordinate neighborhoods. This construction
ensures a precise connectivity property: vertices tending to the same
boundary point can be joined by paths whose vertices tend uniformly
to that point. This is a property of paths in the graph; local
connectivity of $B$ will require a separate argument.

The square-summable edge lengths in $X$ control translated level sets.
They give a \emph{convergence action} on $B$: every sequence of distinct
group elements has a subsequence converging uniformly on compact sets
away from one boundary point to a constant. They also give
$\dim B\le1$ for covering dimension. The group-ring cohomology vanishing supplied by
Poincar\'e duality also gives $\check H^1(B;\mathbb F_2)=0$.
A \emph{cut point} is a point whose removal disconnects the space.
If $B$ has no cut points, set $L=B$. Otherwise, Bowditch's theory of
pretrees reduces the problem to a subcontinuum $L\subset B$ without
cut points. It is a retract: there is a continuous map $B\to L$
fixing every point of $L$. We give the required extension of the
finite-interval reduction to parabolic cut points explicitly.
The resulting $L$ remains nondegenerate and satisfies
\[
 \dim L\le1,\qquad \check H^1(L;\mathbb F_2)=0.
\]

The remaining step is local connectivity of $L$. We measure how deeply
graph vertices lie in a boundary neighborhood by chains of nested
annuli. Cohomology makes the nearby component labels well defined,
and convergence dynamics makes those labels stable as the depth
increases. A key point is uniformity: after fixing the neighborhood and
a finite list of edge and stabilizer moves, one depth threshold works
for every deeper graph path, independently of its endpoints and length.
We obtain this threshold from three fixed boundary points and a
majority-label argument. Thus paths can be pushed to arbitrary depth while their
endpoints approach two prescribed nearby boundary points. Limits of
these paths give small connected subsets of $L$. This argument applies
even when the auxiliary graph has infinite stabilizers and infinite
valence; it uses compactness of the ambient space in place of local
finiteness of that graph.

Finally, the dimension and cohomology properties imply that $L$ contains
no embedded circle. A locally connected compact metric continuum without an embedded
circle is called a \emph{dendrite}. Every nondegenerate dendrite has a
cut point, giving the contradiction.
The finite-energy compactification and the annular argument are the
two constructions that connect the cohomological input to this
continuum-theoretic conclusion.

The topological input is used before any boundary is constructed:
a finite four-dimensional homotopy model supplies finite cellular chains,
local-coefficient duality supplies the group-ring cohomology vanishings,
and von Neumann dimension identifies the positive first $L^2$ class
with a finite-energy
gradient. Later, Bowditch's pretree theorems apply to a continuum before
local connectivity is known. The latter is obtained only after the
retraction and uniform path argument. This order is essential, since
the block stabilizer may be infinitely generated and its graph may
have infinite valence.

Section~\ref{sec:foundations} establishes the algebraic and topological
inputs. Sections~\ref{sec:compactification}--\ref{sec:dynamics} construct
the boundary and its dynamics; Section~\ref{sec:boundary-topology}
proves its dimension and cohomology properties.
Sections~\ref{sec:pretree}--\ref{sec:blocks} remove cut points.
Sections~\ref{sec:annuli}--\ref{sec:path-pushing} prove local connectivity,
Section~\ref{sec:conclusion} completes the proof, and
Section~\ref{sec:consequences} derives the consequences stated above.
